\chapter{Modele symetryczne i niezależność aksjomatu wyboru}
\label{ch:symmetric-choice}
\chaptermark{Modele symetryczne i aksjomat wyboru}

\section{Dlaczego nie wystarcza zwykłe rozszerzenie generyczne?}

W Rozdziałach~\ref{ch:forcing-foundations}--\ref{ch:forcing-theorem}
zbudowaliśmy rozszerzenia generyczne, a w
Rozdziale~\ref{ch:boolean-valued-independence}
wyprowadziliśmy względną niesprzeczność CH i jej zaprzeczenia
z samej niesprzeczności ZFC.
W zwykłym rozszerzeniu modelu ZFC przez zbiorowe pojęcie
wymuszania aksjomat wyboru pozostaje prawdziwy.
Aby skonstruować model ZF, w którym wybór zawodzi,
trzeba zatem ograniczyć zbiór dopuszczalnych nazw.

Celem niniejszego rozdziału jest dowód implikacji
\[
 \operatorname{Con}(\mathrm{ZFC})\Longrightarrow
 \operatorname{Con}(\mathrm{ZF}+\neg\mathrm{AC})
\]
bez zakładania istnienia przeliczalnego tranzytywnego modelu
ZFC. Po połączeniu z uniwersum konstruowalnym
otrzymamy niezależność AC od ZF w zwykłym
znaczeniu względnej niesprzeczności.
Zasadniczy przykład będzie miał \emph{nieskończony zbiór
bez przeliczalnie nieskończonego podzbioru}.

\begin{uwaga}[Dwa poziomy konstrukcji]
\label{rem:sym36-two-levels}
Najpierw opisujemy budowę w przeliczalnym
tranzytywnym modelu $M$ dla przejrzystości notacji
$\tau^G$. Na końcu formułujemy ten sam argument
wewnątrz dowolnego modelu zbiorowego ZFC,
posługując się ilorazem nazw przez filtr generyczny.
Samo istnienie modelu tranzytywnego jest
założeniem silniejszym niż $\operatorname{Con}(\mathrm{ZFC})$;
nie będziemy go używać w końcowym wniosku.
\end{uwaga}

\section{Rodzina Cohenów i grupa permutacji współrzędnych}

\begin{definicja}[Wymuszanie przeliczalnej rodziny ciągów]
\label{def:sym36-forcing}
W modelu ZFC niech
\[
 \mathbb P=\operatorname{Fn}
       (\omega\times\omega,\{0,1\},<\omega).
\]
Warunkami są funkcje częściowe o \emph{skończonych}
dziedzinach, a $q\le p$ oznacza $q\supseteq p$.
Elementem największym jest funkcja pusta $1_{\mathbb P}$.
Jeżeli $G\subseteq\mathbb P$ jest filtrem generycznym,
to $c=\bigcup G$ jest funkcją
$\omega\times\omega\to\{0,1\}$.
Jej $n$-ty wiersz oznaczamy
$c_n(k)=c(n,k)$.
\end{definicja}

\begin{proof}[Uzasadnienie całkowitości funkcji]
Dla każdej pary $(n,k)$ zbiór
$D_{n,k}=\{p:(n,k)\in\operatorname{dom}p\}$
jest gęsty: każdy warunek można przedłużyć,
przypisując brakującej parze dowolny bit.
Generyczność daje przecięcie $G\cap D_{n,k}$.
Zgodność warunków w filtrze pokazuje,
że suma funkcji z $G$ jest jednoznaczna.
\end{proof}

\begin{definicja}[Działanie grupy i skończony nośnik]
\label{def:sym36-support}
Niech $\Gamma$ będzie grupą wszystkich permutacji
zbioru $\omega$ należących do modelu wyjściowego.
Dla $\pi\in\Gamma$ i $p\in\mathbb P$ określamy
\[
 (\pi p)(\pi(n),k)=p(n,k).
\]
Działanie na nazwach wymuszania definiujemy
rekurencyjnie przez
$\pi\tau=\{(\pi\sigma,\pi p):(\sigma,p)\in\tau\}$.
Dla skończonego $E\subseteq\omega$ oznaczamy
\[
 \operatorname{Fix}(E)
 =\{\pi\in\Gamma:\forall n\in E\ \pi(n)=n\}.
\]
Mówimy, że $E$ jest \emph{nośnikiem} nazwy $\tau$,
gdy $\pi\tau=\tau$ dla wszystkich
$\pi\in\operatorname{Fix}(E)$.
Równość tutaj jest \emph{dosłowną równością nazw},
nie tylko równością ich wartości po interpretacji.
\end{definicja}

\begin{lemat}[Podstawowe własności nośników]
\label{lem:sym36-supports}
Dla dowolnych nazw $\sigma,\tau$, skończonych
$E,F\subseteq\omega$ oraz $\pi\in\Gamma$ zachodzą:
\begin{enumerate}
\item jeśli $E$ jest nośnikiem $\tau$, to każdy
skończony nadzbiór $E$ również jest jej nośnikiem;
\item jeśli $E$ jest nośnikiem $\tau$, a $F$ nośnikiem
$\sigma$, to $E\cup F$ jest wspólnym nośnikiem obu nazw;
\item $\pi\operatorname{Fix}(E)\pi^{-1}
=\operatorname{Fix}(\pi[E])$;
\item jeżeli $\pi\tau=\tau$, to
$\pi[\operatorname{dom}\tau]=\operatorname{dom}\tau$.
\end{enumerate}
\end{lemat}
\begin{proof}
Dla pierwszego punktu wystarczy zauważyć,
że $\operatorname{Fix}(F)\subseteq\operatorname{Fix}(E)$,
gdy $E\subseteq F$. Drugi punkt wynika z zastosowania
tej obserwacji do obu nazw. W trzecim punkcie
permutacja $\rho$ ustala $\pi[E]$ dokładnie wtedy,
gdy $\pi^{-1}\rho\pi$ ustala $E$.
Ostatnia własność wynika z porównania
pierwszych współrzędnych par tworzących
nazwę $\tau$ oraz $\pi\tau$.
\end{proof}

\begin{definicja}[Nazwy hereditarne symetryczne]
\label{def:sym36-hs}
Nazwa $\tau$ jest \emph{symetryczna}, gdy ma skończony
nośnik. Jest \emph{heredytarnie symetryczna},
gdy sama ma skończony nośnik i każda
$\sigma\in\operatorname{dom}\tau$ jest
heredytarnie symetryczna.
Klasę wszystkich takich nazw oznaczamy
$\mathrm{HS}$. Dla filtru generycznego $G$
określamy
\[
 N=\mathrm{HS}^G
   =\{\tau^G:\tau\in\mathrm{HS}\}.
\]
\end{definicja}

\begin{lemat}[Definiowalność i niezmienniczość klasy HS]
\label{lem:sym36-hs-invariant}
Klasa $\mathrm{HS}$ jest definiowalna w ZFC.
Dla każdej $\pi\in\Gamma$ zachodzi
$\pi[\mathrm{HS}]=\mathrm{HS}$.
Jeżeli $\tau\in\mathrm{HS}$ ma nośnik $E$,
to $\pi\tau$ ma nośnik $\pi[E]$.
Każda nazwa kanoniczna $\check x$
dla zbioru $x$ z modelu wyjściowego
należy do $\mathrm{HS}$ i ma pusty nośnik.
\end{lemat}
\begin{proof}
Predykat „$E$ jest skończonym nośnikiem $\tau$”
jest formułą teorii mnogości, ponieważ
działanie na nazwach jest zadane
rekurencją po ich rangach.
Rekurencja po randze definiuje także predykat
„$\tau$ oraz wszystkie jej podnazwy
mają skończony nośnik”.
W szczególności nie trzeba kwantyfikować
po zbiorze wszystkich nazw.

Dla $\pi\in\Gamma$ i $\rho\in\operatorname{Fix}(\pi[E])$
mamy $\pi^{-1}\rho\pi\in\operatorname{Fix}(E)$.
Zatem $\rho(\pi\tau)=\pi\tau$.
Indukcja po randze pokazuje równocześnie,
że wszystkie podnazwy $\pi\tau$ należą do HS.
Działając przez $\pi^{-1}$ dostajemy
odwrotną inkluzję klas.

Dla nazwy kanonicznej
$\check x=\{(\check y,1_{\mathbb P}):y\in x\}$
indukcja po randze $x$ daje
$\pi\check x=\check x$, ponieważ
$\pi1_{\mathbb P}=1_{\mathbb P}$.
Pusty nośnik działa również na wszystkie podnazwy.
\end{proof}

\begin{lemat}[Symetria wymuszania]
\label{lem:sym36-symmetry}
Niech $\pi\in\Gamma$, $p\in\mathbb P$,
$\varphi(x_1,\ldots,x_m)$ będzie ustaloną formułą
języka $\{\in\}$, a $\tau_1,\ldots,\tau_m$
nazwami wymuszania. Wówczas
\[
 p\Vdash_{\mathbb P}\varphi(\vec\tau)
 \quad\Longleftrightarrow\quad
 \pi p\Vdash_{\mathbb P}\varphi(\pi\vec\tau).
\]
\end{lemat}
\begin{proof}
Odwzorowanie $p\mapsto\pi p$
jest automorfizmem porządku $\mathbb P$:
zachowuje porządek, zgodność i wszystkie
zbiory gęste przenosi na zbiory gęste.
W atomowych klauzulach wymuszania
z Rozdziału~\ref{ch:forcing-theorem}
zastępujemy każdą parę $(\sigma,p)$
parą $(\pi\sigma,\pi p)$.
Równoczesna indukcja po rangach nazw
daje równoważność dla $\in$ i $=$.
Indukcja po budowie ustalonej formuły
zachowuje negację i koniunkcję,
ponieważ automorfizm zachowuje wszystkie
wzmocnienia warunku, a egzystencję,
ponieważ przenosi nazwy i gęstość
świadków na odpowiadające im nazwy i gęstość.
Zastosowanie $\pi^{-1}$ daje oba kierunki.
\end{proof}

\section{Wymuszanie ograniczone do nazw symetrycznych}

\begin{definicja}[Relacja wymuszania dla HS]
\label{def:sym36-relative-forcing}
Dla \emph{ustalonej standardowej formuły}
$\varphi$ oraz nazw z $\mathrm{HS}$
definiujemy $p\Vdash_{\mathrm{HS}}\varphi$
rekurencyjnie:
\begin{enumerate}
\item dla formuły atomowej stosujemy zwykłą
relację $p\Vdash_{\mathbb P}\varphi$;
\item $p\Vdash_{\mathrm{HS}}(\varphi\land\psi)$,
gdy wymusza osobno oba człony;
\item $p\Vdash_{\mathrm{HS}}\neg\varphi$
wtedy i tylko wtedy, gdy
\emph{żaden} $q\le p$ nie wymusza
$\varphi$ w sensie HS;
\item $p\Vdash_{\mathrm{HS}}\exists x\,\varphi(x,\vec\tau)$
wtedy i tylko wtedy, gdy
\[
 \forall q\le p\ \exists r\le q\
 \exists\sigma\in\mathrm{HS}\
 r\Vdash_{\mathrm{HS}}\varphi(\sigma,\vec\tau).
\]
\end{enumerate}
Pozostałe spójniki i kwantyfikator ogólny
są skrótami zbudowanymi z podanych operacji.
\end{definicja}

\begin{lemat}[Definiowalność, monotoniczność i symetria]
\label{lem:sym36-relative-properties}
Niech $\varphi(\vec x)$ będzie ustaloną
standardową formułą, a $\vec\tau\in\mathrm{HS}$.
W ZFC predykat
$p\Vdash_{\mathrm{HS}}\varphi(\vec\tau)$
jest definiowalny. Jeżeli $q\le p$ i
$p\Vdash_{\mathrm{HS}}\varphi(\vec\tau)$,
to $q\Vdash_{\mathrm{HS}}\varphi(\vec\tau)$.
Ponadto dla $\pi\in\Gamma$ zachodzi
\[
 p\Vdash_{\mathrm{HS}}\varphi(\vec\tau)
 \iff
 \pi p\Vdash_{\mathrm{HS}}\varphi(\pi\vec\tau).
\]
\end{lemat}
\begin{proof}
Dla atomów korzystamy z
Twierdzenia~\ref{thm:forcing33-definability}
i Lematu~\ref{lem:sym36-symmetry}.
Dla formuły złożonej indukcja po jej budowie
zamienia kwantyfikator „$\sigma\in\mathrm{HS}$”
na zwykły kwantyfikator po zbiorach,
z dodatkiem definiowalnego predykatu HS.
Każdy krok jest zatem formułą ZFC,
choć klasa HS nie jest zbiorem.
Monotoniczność wynika ze zmniejszania
zbioru dopuszczalnych wzmocnień.
Przy egzystencji, jeśli warunek $p$
spełnia kryterium gęstości,
to spełnia je każdy $q\le p$.
Niezmienniczość dla negacji i egzystencji
wynika z bijektywnego przenoszenia
wzmocnień i z Lematu~\ref{lem:sym36-hs-invariant}.
\end{proof}

\begin{twierdzenie}[Lemat prawdy dla modeli symetrycznych]
\label{thm:sym36-truth}
Niech $M$ będzie przeliczalnym tranzytywnym
modelem ZFC, $\mathbb P,\mathrm{HS}$ konstrukcjami
z Definicji~\ref{def:sym36-forcing} oraz
Definicji~\ref{def:sym36-hs} wykonanymi w $M$,
a $G\subseteq\mathbb P^M$ filtrem $M$-generycznym.
Niech $N=\mathrm{HS}^{M,G}$.
Dla każdej ustalonej formuły
$\varphi(x_1,\ldots,x_m)$ i
$\tau_1,\ldots,\tau_m\in(\mathrm{HS})^M$ mamy
\[
 N\models\varphi(\tau_1^G,\ldots,\tau_m^G)
 \iff
 \exists p\in G\
 \bigl(M\models
   p\Vdash_{\mathrm{HS}}\varphi(\vec\tau)\bigr).
\]
\end{twierdzenie}
\begin{proof}
Prowadzimy indukcję po budowie formuły.
Dla formuł atomowych zwykły lemat prawdy
z Rozdziału~\ref{ch:forcing-theorem}
obowiązuje bez zmian: prawdziwość
$\sigma^G\in\tau^G$ lub $\sigma^G=\tau^G$
nie zależy od tego, czy nazwę rozpatrujemy
w $N$, czy w pełnym $M[G]$.

Dla koniunkcji wybieramy wspólne
wzmocnienie dwóch warunków z $G$.
Dla negacji zakładamy najpierw,
że $N\models\neg\psi(\vec\tau)$.
Zbiór warunków, które wymuszają
$\neg\psi$ albo mają wzmocnienie
wymuszające $\psi$, jest gęsty i
należy do $M$ z definiowalności
Lematu~\ref{lem:sym36-relative-properties}.
Generyczność daje warunek z $G$
rozstrzygający jedną z dwóch możliwości;
druga przeczy hipotezie indukcyjnej.
Odwrotnie, gdy $p\in G$ wymusza negację,
żaden warunek z $G$ nie może wymuszać $\psi$:
istniałoby ich wspólne wzmocnienie w $G$.

Jeżeli $N\models\exists x\,\psi(x,\vec\tau)$,
to świadkiem jest $\sigma^G$ dla
pewnej $\sigma\in(\mathrm{HS})^M$.
Indukcja daje $p\in G$ wymuszający
$\psi(\sigma,\vec\tau)$.
Warunki wymuszające tę konkretną
formułę są gęste poniżej $p$
w trywialny sposób (sam $p$
jest odpowiednim wzmocnieniem);
zatem $p$ wymusza egzystencję.
W drugą stronę, jeśli $p\in G$
wymusza egzystencję, to zbiór
wzmocnień $r$ poniżej $p$
mających świadka $\sigma\in\mathrm{HS}$
jest gęsty poniżej $p$ i należy do $M$.
Generyczność $G$ daje taki $r\in G$.
Hipoteza indukcyjna świadczy,
że $N\models\psi(\sigma^G,\vec\tau^G)$.
\end{proof}

\section{Zachowanie aksjomatów ZF}

\begin{twierdzenie}[Twierdzenie o modelu symetrycznym]
\label{thm:sym36-zf}
Przy założeniach
Twierdzenia~\ref{thm:sym36-truth}
klasa $N=\mathrm{HS}^{M,G}$ jest
tranzytywnym modelem wewnętrznym ZF
zawartym w $M[G]$, zawierającym $M$
i mającym te same liczby porządkowe co $M[G]$.
W szczególności $N\models\mathrm{ZF}$.
\end{twierdzenie}
\begin{proof}
Z Lematu~\ref{lem:sym36-hs-invariant}
wynika, że każda nazwa kanoniczna jest w HS,
więc $M\subseteq N$. Jeżeli $x\in y=\tau^G$
dla $\tau\in\mathrm{HS}$, to
$x=\sigma^G$ dla pewnej pary
$(\sigma,p)\in\tau$ z $p\in G$.
Heredytarna symetria daje $\sigma\in\mathrm{HS}$,
czyli $x\in N$; stąd tranzytywność.
Żadne wymuszanie zbiorowe nie dodaje
liczb porządkowych, a wszystkie
liczby porządkowe modelu wyjściowego
pozostają w $N$. Zatem klasa $N$ ma
te same liczby porządkowe co $M[G]$.

Sprawdzimy teraz aksjomaty.
W całym dowodzie „$p$ wymusza”
oznacza relację $\Vdash_{\mathrm{HS}}$.
W szczególności wszystkie nazwy parametrów
należą do HS, a używane
formuły są ustalone i standardowe.

\emph{Ekstensjonalność i ufundowanie.}
Ponieważ $N$ jest klasą tranzytywną
w modelu $M[G]$, oba aksjomaty dziedziczy
po $M[G]$. Dla ufundowania trzeba jedynie
zauważyć, że jeśli $a\in N$ jest niepusty,
to dowolny element $\in$-minimalny
zbioru $a$, znaleziony w $M[G]$,
należy do $N$ dzięki tranzytywności.

\emph{Pusty zbiór, para, suma i nieskończoność.}
Pustą nazwę i $\check\omega$ mamy w HS.
Jeśli $\sigma,\tau\in\mathrm{HS}$
mają nośniki $E,F$, to nazwa pary
$\{(\sigma,1_{\mathbb P}),(\tau,1_{\mathbb P})\}$
ma nośnik $E\cup F$ i wszystkie podnazwy
w HS. Dla sumy nazwy $\sigma$
używamy zwykłej konstrukcji
$\bigcup\sigma$ w postaci nazwy
\[
 \mathrm U(\sigma)=
 \{(\rho,r):\exists(\pi,p)\in\sigma\
       \exists(\rho,q)\in\pi\
       (r\le p\ \land\ r\le q)\}.
\]
Wszystkie jej podnazwy leżą w HS,
a każdy automorfizm ustalający
$\sigma$ ustala także $\mathrm U(\sigma)$.
Jest to nazwa sumy zbiorów
zakodowanych przez elementy $\sigma^G$,
ponieważ filtr generyczny przecina
gęste zbiory wzmocnień rozstrzygających
obie przynależności. To daje aksjomat sumy.

\emph{Wyróżnianie.}
Dla $\tau,\vec\eta\in\mathrm{HS}$
i ustalonej formuły
$\varphi(x,\vec\eta)$ definiujemy
\[
 S_{\varphi,\tau}=
 \{(\rho,p):\rho\in\operatorname{dom}\tau,\
   p\Vdash_{\mathrm{HS}}
    (\rho\in\tau\land\varphi(\rho,\vec\eta))\}.
\]
Zbiór ten istnieje w $M$,
ponieważ $\operatorname{dom}\tau\times\mathbb P$
jest zbiorem i relacja wymuszania
jest definiowalna. Wszystkie jego podnazwy
są w HS, a wspólny nośnik nazw parametrów
ustala $S_{\varphi,\tau}$ na mocy symetrii
relacji wymuszania. Lemat prawdy pokazuje
dla każdego filtru generycznego, że
$S_{\varphi,\tau}^G=
 \{x\in\tau^G:N\models\varphi(x,\vec\eta^G)\}$.
To dowodzi wszystkich instancji
schematu wyróżniania.

\emph{Zbiór potęgowy.}
Niech $D=\operatorname{dom}\tau$ i
\[
 T_\tau=
 \{\sigma\in\mathrm{HS}:
       \sigma\subseteq D\times\mathbb P\}.
\]
To jest \emph{zbiór} w $M$,
gdyż jest podzbiorem zbioru
potęgowego $D\times\mathbb P$.
Definiujemy nazwę
\[
 \dot{\mathcal P}_N(\tau)=
 \{(\sigma,p):\sigma\in T_\tau,\
              p\Vdash_{\mathrm{HS}}
                  \sigma\subseteq\tau\}.
\]
Jeśli $E$ jest nośnikiem $\tau$,
to $\pi[D]=D$ dla $\pi\in\operatorname{Fix}(E)$,
więc $T_\tau$ jest niezmienniczy
i również $\dot{\mathcal P}_N(\tau)$
ma nośnik $E$. Każda nazwa w jej
dziedzinie jest w HS, dlatego
$\dot{\mathcal P}_N(\tau)\in\mathrm{HS}$.
Każdy jej element jest podzbiorem $\tau^G$.

Odwrotnie, niech $x=\sigma^G\in N$
i $x\subseteq\tau^G$.
Definiujemy \emph{ograniczoną nazwę}
\[
 \sigma_*=
 \{(\rho,p):\rho\in D,\
   p\Vdash_{\mathrm{HS}}
    (\rho\in\sigma\land\rho\in\tau)\}.
\]
Nazwę tę ustala każdy automorfizm
ustalający obie nazwy $\sigma,\tau$,
a jej podnazwy należą do HS.
Stąd $\sigma_*\in T_\tau$.
Z lematu prawdy i z faktu,
że wszystkie elementy $\tau^G$
są interpretacjami nazw z $D$,
otrzymujemy
$\sigma_*^G=\sigma^G\cap\tau^G=x$.
Co więcej, już
$1_{\mathbb P}\Vdash_{\mathrm{HS}}
 \sigma_*\subseteq\tau$,
bo każdy warunek występujący w
$\sigma_*$ wymusza przynależność do $\tau$.
Zatem para
$(\sigma_*,1_{\mathbb P})$ należy do
$\dot{\mathcal P}_N(\tau)$.
Każdy podzbiór $\tau^G$ należący
do $N$ występuje więc w tej samej
nazwie zbiorowej. Jest to wymagany
zbiór potęgowy \emph{obliczony wewnątrz $N$};
nie należy go mylić z całym
$\mathcal P(\tau^G)$ w $M[G]$.

\emph{Schemat zbierania (Collection).}
Niech $\varphi(x,y,\vec\eta)$ będzie
ustaloną formułą, $\tau,\vec\eta\in\mathrm{HS}$.
W $M$ rozpatrujemy \emph{zbiór} par
$(p,\rho)$, gdzie $p\in\mathbb P$,
$\rho\in\operatorname{dom}\tau$ i
\[
 p\Vdash_{\mathrm{HS}}
   \bigl(\rho\in\tau\land
          \exists y\,\varphi(\rho,y,\vec\eta)\bigr).
\]
Dla każdej takiej pary relacja
wymuszania egzystencji zapewnia,
że poniżej $p$ istnieje warunek $q$
i nazwa $\sigma\in\mathrm{HS}$ taka, że
$q\Vdash_{\mathrm{HS}}\varphi(\rho,\sigma,\vec\eta)$.
Wybieramy \emph{najmniejszą rangę}
możliwej nazwy $\sigma$.
Dla ustalonej pary minimum można
znaleźć, wybierając jednego świadka
i minimalizując w zbiorze liczb
porządkowych nie większych od jego rangi.
Zastępowanie w $M$ ogranicza te rangi
jedną liczbą porządkową $\alpha$,
tak że każdy rozpatrywany $(p,\rho)$
ma jakieś $q\le p$ i świadka
$\sigma\in\mathrm{HS}$ o randze $<\alpha$.

Niech $V_\alpha^{\mathbb P}$ oznacza
zbiór nazw wymuszania o randze $<\alpha$
(zgodnie z rekurencją nazw
z Rozdziału~\ref{ch:forcing-foundations}).
Definiujemy
\[
 \dot C=
 \{(\sigma,p):
   \sigma\in \mathrm{HS}\cap V_\alpha^{\mathbb P},\
   p\Vdash_{\mathrm{HS}}
      \exists x\in\tau\
        \varphi(x,\sigma,\vec\eta)\}.
\]
Jest to nazwa zbiorowa w $M$.
Wspólny nośnik $\tau,\vec\eta$ ustala
jej dziedzinę i warunek przynależności,
bo automorfizmy zachowują rangi;
podnazwy należą do HS.
Zatem $\dot C\in\mathrm{HS}$.

Załóżmy teraz w $N$, że
$\forall x\in\tau^G\ \exists y\,\varphi(x,y,\vec\eta^G)$.
Dla ustalonego $x=\rho^G\in\tau^G$
lemat prawdy dostarcza $p\in G$
wymuszający jednocześnie
$\rho\in\tau$ i istnienie obrazu.
Dla \emph{każdego} $r\le p$
istnieje $q\le r$ i świadek
$\sigma$ rangi $<\alpha$, ponieważ
$r$ wymusza te same dwa warunki
i jest jedną z par uwzględnionych
przy ograniczaniu rang.
Zbiór takich $q$ jest gęsty poniżej $p$
i należy do $M$. Filtr $G$ przecina
ten zbiór w pewnym $q$.
Wtedy $q$ wymusza równocześnie
$\rho\in\tau$ i $\varphi(\rho,\sigma,\vec\eta)$,
zatem $(\sigma,q)\in\dot C$.
W konsekwencji $\sigma^G\in\dot C^G$
jest wymaganym świadkiem dla $x$.
Ponieważ każdy element $\dot C^G$
z definicji jest wartością obrazu
jakiegoś $x\in\tau^G$, otrzymujemy
schemat zbierania. Po zastosowaniu
wyróżniania daje on także
schemat zastępowania.

Sprawdziliśmy zatem ekstensjonalność,
pusty zbiór, parę, sumę, nieskończoność,
zbiór potęgowy, wyróżnianie,
zastępowanie i ufundowanie.
Nie używaliśmy aksjomatu wyboru
\emph{wewnątrz} $N$.
\end{proof}

\section{Nieskończony zbiór Dedekinda skończony}

\begin{definicja}[Zbiór Dedekinda skończony]
\label{def:sym36-dedekind}
Zbiór $A$ jest \emph{Dedekinda skończony},
jeśli nie istnieje iniekcja $f:\omega\to A$.
W ZF zbiór może być nieskończony
(nie być równoliczny z żadną liczbą naturalną)
i jednocześnie Dedekinda skończony.
\end{definicja}

\begin{lemat}[Nazwy wierszy i całej rodziny]
\label{lem:sym36-cohen-family}
Dla każdego $n\in\omega$ niech
\[
 \dot c_n=
 \{(\check k,p):
     (n,k)\in\operatorname{dom}p,\
     p(n,k)=1\},\qquad
 \dot A=\{(\dot c_n,1_{\mathbb P}):n\in\omega\}.
\]
Wtedy $\dot c_n\in\mathrm{HS}$
ma nośnik $\{n\}$, a
$\dot A\in\mathrm{HS}$ ma nośnik pusty.
Dla każdego filtru generycznego $G$
nazwy te reprezentują odpowiednio
podzbiory $\omega$ oraz zbiór
$A=\{c_n:n<\omega\}$;
elementy $c_n$ są parami różne.
\end{lemat}
\begin{proof}
Permutacja ustalająca $n$ nie zmienia
pierwszej współrzędnej $(n,k)$, więc
ustala $\dot c_n$. Wszystkie jej podnazwy
$\check k$ mają pusty nośnik.
Ponadto $\pi\dot c_n=\dot c_{\pi(n)}$,
wobec czego działanie dowolnej $\pi$
jedynie przestawia elementy nazwy $\dot A$.
Zatem $\dot A$ ma pusty nośnik,
a jej podnazwy są w HS.

Dla ustalonych $n\ne m$ zbiór
\[
 D_{n,m}=\{p:\exists k\
 [(n,k),(m,k)\in\operatorname{dom}p
   \land p(n,k)\ne p(m,k)]\}
\]
jest gęsty. Dla danego warunku $p$
wybieramy $k$ spoza obu skończonych
zbiorów użytych współrzędnych
i przedłużamy $p$, dopisując
odpowiednio wartości $0$ i $1$.
Generyczność zapewnia, że wiersze
$c_n,c_m$ różnią się na pewnym $k$.
Definicja nazwy $\dot c_n$ i gęstość
warunków rozstrzygających każdy bit
pokazują, że $\dot c_n^G
=\{k:c(n,k)=1\}$. Stąd teza.
\end{proof}

\begin{lemat}[Skończony nośnik ogranicza wartości funkcji]
\label{lem:sym36-finite-range}
Niech $\dot f\in\mathrm{HS}$ ma skończony
nośnik $E\subseteq\omega$ i niech
$p\in\mathbb P$ wymusza w sensie HS,
że $\dot f$ jest funkcją o dziedzinie
$\check\omega$ i wartościach w $\dot A$.
Wówczas $p$ wymusza, że zbiór wartości
$\dot f$ jest podzbiorem
$\{\dot c_k:k\in E\}$.
\end{lemat}
\begin{proof}
Ustalmy $m\in\omega$.
Z definicji \emph{funkcji o wartościach w $\dot A$},
z przynależności do nazwy
$\dot A$ i z gęstości świadków wynika,
że poniżej każdego $q\le p$ istnieje
$r\le q$ i $k\in\omega$ takie, że
\[
 r\Vdash_{\mathrm{HS}}
 \dot f(\check m)=\dot c_k.
\]
Pokażemy, że koniecznie $k\in E$.

Załóżmy przeciwnie, że $k\notin E$.
Ponieważ $r$ jest skończoną funkcją,
istnieje $\ell\notin E\cup\{k\}$
niewystępujące jako pierwsza
współrzędna w $\operatorname{dom}r$.
Niech $\pi$ zamienia $k$ z $\ell$
i ustala pozostałe liczby naturalne.
Permutacja ta ustala $E$, a więc
$\pi\dot f=\dot f$ oraz
$\pi\check m=\check m$.
Lemat~\ref{lem:sym36-relative-properties}
daje
\[
 \pi r\Vdash_{\mathrm{HS}}
  \dot f(\check m)=\dot c_\ell.
\]
Warunki $r$ i $\pi r$ są zgodne:
wszystkie współrzędne różne od
wierszy $k,\ell$ są identyczne,
a na wierszu $\ell$ warunek $r$
nie ma żadnego przypisania.
Mają zatem wspólne wzmocnienie $s$.

Warunek $s$ wymusza jednocześnie
$\dot f(\check m)=\dot c_k$
i $\dot f(\check m)=\dot c_\ell$,
czyli $\dot c_k=\dot c_\ell$.
Przeczy to Lematu~\ref{lem:sym36-cohen-family},
gdyż $k\ne\ell$ i już
$1_{\mathbb P}$ wymusza ich nierówność.
Stąd $k\in E$.

Dla każdego $m$ warunki, które
rozstrzygają $\dot f(\check m)$
jako jeden z wierszy o indeksach w $E$,
są gęste poniżej $p$.
Definicja wymuszania kwantyfikatora
ogólnego i dowiedziony lemat prawdy
dają zapowiedzianą inkluzję
zbioru wartości dla całej funkcji.
\end{proof}

\begin{twierdzenie}[Nieskończony zbiór Dedekinda skończony]
\label{thm:sym36-dedekind}
Przy założeniach
Twierdzenia~\ref{thm:sym36-zf}
zbiór $A=\dot A^G\in N$ jest
nieskończony, lecz w $N$
nie istnieje iniekcja $\omega\to A$.
W szczególności $N\models\neg\mathrm{AC}$.
\end{twierdzenie}
\begin{proof}
Lemat~\ref{lem:sym36-cohen-family}
pokazuje, że $A$ ma przynajmniej
$n$ elementów dla każdego
$n\in\omega$: wystarczy
$c_0,\ldots,c_{n-1}$.
Są one elementami $N$,
bo ich nazwy należą do HS.
Zatem $A$ nie jest zbiorem skończonym.

Przypuśćmy, że w $N$ istnieje
iniekcja $f:\omega\to A$.
Z definicji $N$ wybieramy
$\dot f\in\mathrm{HS}$ z
$\dot f^G=f$ i skończony
nośnik $E$ nazwy $\dot f$.
Lemat prawdy dostarcza $p\in G$
wymuszający, że $\dot f$
jest funkcją $\omega\to A$
i jest iniekcją.
Lemat~\ref{lem:sym36-finite-range}
ogranicza jej zbiór wartości
do skończonego zbioru
$\{c_k:k\in E\}$.
Nie istnieje iniekcja zbioru
$\omega$ w zbiór skończony;
sprzeczność.

W ZF z aksjomatu wyboru wynika,
że każdy nieskończony zbiór
ma podzbiór przeliczalnie nieskończony:
dobrze porządkujemy zbiór i wybieramy
rekurencyjnie pierwszy element,
który nie znalazł się w skończonej
liście wcześniej wybranych elementów.
Otrzymana funkcja z $\omega$ jest
iniekcją. Ponieważ nasz zbiór $A$
jest kontrprzykładem, AC zawodzi w $N$.
\end{proof}

\begin{przyklad}[Dlaczego nie można „po prostu ponumerować” wierszy?]
\label{ex:sym36-indexing}
W pełnym rozszerzeniu $M[G]$
istnieje funkcja $n\mapsto c_n$,
zadana generyczną tablicą bitów.
W modelu $N$ każdy pojedynczy $c_n$
oraz zbiór $A$ istnieją, lecz
funkcja wyliczająca wszystkie $c_n$
nie jest elementem $N$.
Jej zwykła nazwa nie ma skończonego
nośnika: permutacja zmieniająca
indeks spoza skończonego $E$
zmienia odpowiednią parę w grafie.
Ścisłym uzasadnieniem
\emph{nieistnienia jakiejkolwiek}
iniekcji, a nie tylko tej konkretnej
funkcji, jest
Lemat~\ref{lem:sym36-finite-range}.
\end{przyklad}

\section{Od niesprzeczności ZFC do niezależności AC od ZF}

\begin{lemat}[Interpretacja bez tranzytywności modelu wyjściowego]
\label{lem:sym36-arbitrary-model}
Niech $M\models\mathrm{ZFC}$ będzie dowolnym
\emph{przeliczalnym modelem zbiorowym}.
Nie zakładamy zewnętrznego dobrego ufundowania
relacji $\in^M$.
W metateorii ZFC istnieje model zbiorowy
$N\models\mathrm{ZF}+\neg\mathrm{AC}$
zbudowany z wewnętrznych nazw
$(\mathrm{HS})^M$ dla
$(\mathbb P,\Gamma)^M$.
\end{lemat}
\begin{proof}
Ponieważ $|M|$ jest zewnętrznie
przeliczalny, możemy wyliczyć
wszystkie zbiory warunków
$D\subseteq\mathbb P^M$, które
$M$ uważa za gęste, i kolejno
wybierać ich warunki. Używając
skierowania porządku przez wzmocnienia,
konstruujemy \emph{zewnętrzny} filtr
$M$-generyczny $G$ na $\mathbb P^M$.
Nie potrzeba w tym celu zewnętrznej
rekurencji po nieufundowanych rangach nazw.

Dla $\sigma,\tau\in(\mathrm{HS})^M$
definiujemy
\[
 \sigma\equiv_G\tau
 \iff\exists p\in G\
     M\models p\Vdash_{\mathrm{HS}}\sigma=\tau,
\]
oraz
\[
 [\sigma]_G\in_G[\tau]_G
 \iff\exists p\in G\
     M\models p\Vdash_{\mathrm{HS}}\sigma\in\tau.
\]
Zwykła poprawność wymuszania dla
atomów, podstawianie równych nazw
i własności filtru dowodzą,
że pierwsza relacja jest równoważnością,
a druga nie zależy od reprezentantów.
Klasy i relacja są \emph{zbiorami zewnętrznymi},
ponieważ $(\mathrm{HS})^M\subseteq|M|$
jest zewnętrznie zbiorem.
Otrzymujemy więc strukturę zbiorową.

Wszystkie dowody definicji
$\Vdash_{\mathrm{HS}}$, symetrii,
zbierania i potęgowania z
Twierdzenia~\ref{thm:sym36-zf}
są dowodami pierwszego rzędu
\emph{w ZFC modelu wyjściowego}:
kwantyfikują po zbiorowych nazwach,
skończonych nośnikach i porządkach,
a przy ustalonej formule mają
definiowalny predykat wymuszania.
Stąd każdy ich krok pozostaje
prawdziwy, gdy interpretujemy go
w nietranzytywnym $M$.
Indukcja po \emph{standardowej}
formule, analogiczna do
Twierdzenia~\ref{thm:sym36-truth},
daje dla klas nazw lemat prawdy:
\[
 N\models\varphi([\vec\tau]_G)
 \iff\exists p\in G\
       (M\models p\Vdash_{\mathrm{HS}}
          \varphi(\vec\tau)).
\]
Argument o skończonym nośniku
z Lematu~\ref{lem:sym36-finite-range}
również jest dowodem wewnętrznym
ZFC. Otrzymany model spełnia
każdą \emph{standardową}
instancję ZF oraz zaprzeczenie AC.
Jego zewnętrzna relacja przynależności
nie musi być dobrze ufundowana;
do niesprzeczności teorii nie
jest to potrzebne.
\end{proof}

\begin{twierdzenie}[Względna niesprzeczność i niezależność AC]
\label{thm:sym36-relative-consistency}
Dla teorii ZF i ZFC zachodzą implikacje
metamatematyczne
\[
 \operatorname{Con}(\mathrm{ZFC})
 \Longrightarrow
 \operatorname{Con}(\mathrm{ZF}+\neg\mathrm{AC}),
\]
oraz
\[
 \operatorname{Con}(\mathrm{ZF})
 \Longrightarrow
 \begin{cases}
 \operatorname{Con}(\mathrm{ZF}+\mathrm{AC}),\\
 \operatorname{Con}(\mathrm{ZF}+\neg\mathrm{AC}).
 \end{cases}
\]
W szczególności, jeśli ZF jest niesprzeczna,
to ZF nie dowodzi ani AC, ani $\neg$AC.
\end{twierdzenie}
\begin{proof}
Z założenia $\operatorname{Con}(\mathrm{ZFC})$
i twierdzenia Gödla o pełności otrzymujemy
model zbiorowy ZFC. Dolne twierdzenie
Löwenheima--Skolema pozwala wybrać model
przeliczalny. Lemat~\ref{lem:sym36-arbitrary-model}
daje model ZF+$\neg$AC, a więc
niesprzeczność tej teorii.

Z drugiej strony ZF dowodzi, że
uniwersum konstruowalne $L$
spełnia ZFC, co wykazaliśmy w
Rozdziale~\ref{ch:constructible}.
Tym samym
$\operatorname{Con}(\mathrm{ZF})$
implikuje
$\operatorname{Con}(\mathrm{ZFC})
=\operatorname{Con}(\mathrm{ZF}+\mathrm{AC})$.
Składając tę implikację z
pierwszą częścią dowodu, otrzymujemy
niesprzeczność także ZF+$\neg$AC.
Gdyby ZF dowodziła jednego z dwóch
sprzecznych zdań AC, $\neg$AC,
odpowiednia druga teoria byłaby sprzeczna.
\end{proof}

\begin{uwaga}[AC, AD i zakres wniosku]
\label{rem:sym36-ac-ad}
Niezależność AC \emph{od ZF} jest innym
zagadnieniem niż relacja między AC i
aksjomatem determinacji AD.
W Rozdziale~\ref{ch:determinacy}
udowodniliśmy, że ZF+AC wyklucza
pełne AD. Obecna konstrukcja daje
model ZF+$\neg$AC, lecz \emph{nie}
daje automatycznie modelu ZF+AD:
sam brak wyboru nie implikuje determinacji.
Nie utożsamiamy także zbioru Dedekinda
skończonego ze zbiorem skończonym
w zwykłym znaczeniu.
\end{uwaga}

\begin{uwaga}[Granice formalnej weryfikacji]
\label{rem:sym36-formal-status}
Dowody w tym rozdziale są tekstowymi
dowodami matematycznymi, nie wynikiem
kompilacji ich formalizacji w Lean.
Kontrola jednostkowa semantyki nazw,
odnośników i kompilacji LaTeX nie zastępuje
weryfikacji poprawności dowodów przez
jądro asystenta dowodzenia.
\end{uwaga}
